\documentclass[11pt]{article}
\usepackage[a4paper,margin=26mm]{geometry}
\usepackage{amsmath,amssymb,amsthm,booktabs,array}
\usepackage[T1]{fontenc}
\usepackage{lmodern}
\usepackage[colorlinks=true,allcolors=blue,breaklinks=true]{hyperref}
\newtheorem{theorem}{Theorem}
\newtheorem{lemma}[theorem]{Lemma}
\newtheorem{corollary}[theorem]{Corollary}
\newcommand{\R}{\mathbb R}
\newcommand{\C}{\mathbb C}
\newcommand{\Z}{\mathbb Z}
\newcommand{\bop}{Banica--\"Ozteke--Pittau}
\title{Extending Banica--\"Ozteke--Pittau's Fourier defect theory:\protect\\
a correction to their truncation conjecture}
\author{Haobo Ma\\\small ChronoAI Pte Ltd, Singapore
\and Wenlin Zhang\\\small National University of Singapore\\\small ChronoAI Pte Ltd, Singapore}
\date{14 November 2026}
\begin{document}
\maketitle
\begin{abstract}
Extending Banica--\"Ozteke--Pittau's Fourier defect theory, we refute
their Conjecture~1.14
(\emph{Isolated partial Hadamard matrices, and related topics},
arXiv:1706.00986v2, p.~9): Fourier incidence equations and the analytic
implicit-function theorem give non-isolation after any two row deletions
at every prime order $p\geq5$.
For any nonempty proper row set of size $m$, the local
dephased solution space is a real-analytic manifold of dimension
$(m-1)(p-m-1)$. Thus among the nonempty prime Fourier truncations,
precisely those with $m=1,p-1,p$ are isolated, and no fixed positive
density margin ensures isolation.
We also give an explicit defect formula for every row truncation of a
finite abelian Fourier matrix; its zero-cycle case has the same analytic
dimension formula with $p$ replaced by the group order.
\end{abstract}

\section{The question and the sharp replacement}
Put $\zeta=\exp(2\pi i/p)$, where $p$ is prime, and let
$\varnothing\ne S\subseteq\Z/p\Z$, with $m=|S|$.
The matrix $F_{S,p}=(\zeta^{sj})_{s\in S,\,j\in\Z/p\Z}$ has unit-modulus
entries and orthogonal rows, hence is partial Hadamard in the sense of
\cite[Definition 1.1]{BOP}. Indeed, for $s\ne t$,
$\sum_{j=0}^{p-1}\zeta^{(s-t)j}=0$ by the geometric-series identity.
Isolation means isolation after multiplication of rows and columns by
unit scalars; allowing permutations does not alter local isolation,
since they form a finite group. A local phase slice will make this precise.
The density conjecture in \cite[Conjecture 1.14, p.~9]{BOP} is the assertion
\[
\exists\varepsilon>0\ \forall p\text{ prime}\ \forall S\subseteq\Z/p\Z,
\qquad |S|\geq(1-\varepsilon)p\ \Longrightarrow\ F_{S,p}\text{ isolated}.
\]
Its following paragraph reports evidence that
\begin{quote}\small
``$F_{p-1,p}$ should be always isolated'' and
``$F_{p-2,p}$ should be isolated too, provided that $p$ is big enough''.
\end{quote}
The first suggestion holds; the second fails. The following
dimension formula gives a sharp replacement, including the one-row boundary.

The closest prior method is the Fourier-transformed tangent-space
description in \cite[Proposition~1.10 and Theorem~1.13]{BOP}, following
the square-matrix defect theory of Tadej--\.{Z}yczkowski
\cite[Theorems~5.2--5.3]{TZ} and Banica \cite[Theorem~4.2]{Defect}.
The square prime boundary is an established rigidity result:
Tadej--\.{Z}yczkowski credit it to Petrescu's 1997 thesis
\cite[Introduction and Section~5.1(i)]{TZ}.
Here path incidence makes the proper-prime orthogonality map regular;
this integrates the tangent space and determines the actual local
geometry, rather than only its linearized dimension.

\begin{theorem}\label{thm:dimension}
If $p$ is prime and $\varnothing\ne S\subsetneq\Z/p\Z$, then a neighborhood
of $F_{S,p}$ in a fixed dephased slice of partial Hadamard matrices is a
real-analytic manifold of dimension
\begin{equation}\label{eq:dimension}
D(m,p)=(m-1)(p-m-1).
\end{equation}
Consequently, $F_{S,p}$ is non-isolated if $2\leq m\leq p-2$, and is
isolated if $m=1$ or $m=p-1$.
\end{theorem}

\begin{corollary}\label{cor:counterexample}
For every prime $p\geq5$, deleting any two rows of $F_p$ gives a
non-isolated partial Hadamard matrix, with local dephased dimension $p-3$.
There is no $\varepsilon>0$ as in Conjecture 1.14.
\end{corollary}
\begin{proof}
Take $m=p-2$ in Theorem~\ref{thm:dimension}. Given $\varepsilon>0$, choose
a prime $p>\max\{3,2/\varepsilon\}$. Then $p-2>(1-\varepsilon)p$, while
every such truncation is non-isolated.
\end{proof}

\section{The phase differential has full rank}
Near $F_{S,p}$ every unit-modulus matrix has unique small real phase
coordinates
\[
H(\theta)_{sj}=\zeta^{sj}e^{i\theta_{sj}},\qquad\theta\in\R^{S\times\Z/p\Z}.
\]
Choose one ordering of $S$ and put
\[
\Phi_{st}(\theta)=\sum_{j=0}^{p-1}\zeta^{(s-t)j}
       e^{i(\theta_{sj}-\theta_{tj})}\quad(s<t).
\]
This real-analytic map has values in $\C^{\binom m2}$, regarded as a
real space of dimension $m(m-1)$. Its zero set is exactly the partial
Hadamard condition in these coordinates. At zero its differential is
\begin{equation}\label{eq:jacobian}
(J A)_{st}=i\sum_{j=0}^{p-1}\zeta^{(s-t)j}(A_{sj}-A_{tj}).
\end{equation}

\begin{lemma}\label{lem:surjectivity}
For a proper row set $S$, the real-linear map $J$ is surjective.
\end{lemma}
\begin{proof}
The case $p=2$ has at most one retained row and zero-dimensional target.
Suppose $p$ is odd. For a real row $A_s$ write
\[
x_s(d)=\sum_j\zeta^{dj}A_{sj},\qquad
x_s(-d)=\overline{x_s(d)}.
\]
Conversely, arbitrary coefficients satisfying this conjugacy condition
and $x_s(0)\in\R$ give a real row by
$A_{sj}=p^{-1}\sum_d\zeta^{-dj}x_s(d)$.

To prescribe $(JA)_{st}=y_{st}$ for $s<t$, prescribe
$b_{st}=-iy_{st}$ and extend by
$b_{ts}=-\overline{b_{st}}$. We need
\begin{equation}\label{eq:edge}
x_s(s-t)-x_t(s-t)=b_{st}\quad(s\ne t).
\end{equation}
For one representative $d$ from each pair $\{d,-d\}$, consider the graph
on $S$ with edges joining vertices whose difference is $d$ or $-d$.
Addition by $d$ is a single cycle on $\Z/p\Z$ because $p$ is prime.
Since at least one vertex is absent, its induced graph on $S$ is a union
of paths and isolated vertices. For each edge orient it as $(s,t)$ with
$s-t=d$ and prescribe the difference of the complex potentials
$x_s(d)-x_t(d)=b_{st}$. Start each path with potential zero and extend
successively along its edges; there is no cycle consistency condition.
Now set $x_s(-d)=\overline{x_s(d)}$ and $x_s(0)=0$. The reverse equation
in \eqref{eq:edge} becomes
$\overline{x_t(d)}-\overline{x_s(d)}=-\overline{b_{st}}=b_{ts}$,
so all ordered row pairs are satisfied. Fourier inversion gives real
$A$ satisfying \eqref{eq:jacobian}, as required.
\end{proof}

\section{From the differential to actual deformations}
Fix $s_0\in S$ and column $0$. Additive row and column phases form
\[
T=\{(a_s+b_j)_{sj}:a\in\R^S,\ b\in\R^p\}.
\]
The parameter map has the one-dimensional kernel
$a_s=c$, $b_j=-c$, so $\dim T=m+p-1$.
Define the slice
$V=\{\theta:\theta_{s0}=\theta_{s_0j}=0\text{ for all }s,j\}$.
Every phase array has the unique decomposition into an element of $V$
and an element of $T$, obtained by the dephasing formula
\begin{equation}\label{eq:slice}
(\mathcal D\theta)_{sj}
=\theta_{sj}-\theta_{s0}-\theta_{s_0j}+\theta_{s_00}.
\end{equation}
Thus $\dim V=(m-1)(p-1)$. This is a slice relative to $F_{S,p}$;
its fixed base row need not consist of ones.

Adding $a_s+b_j$ to $\theta_{sj}$ multiplies $\Phi_{st}$ by
$e^{i(a_s-a_t)}$, and hence preserves its zero set. Also $T\subseteq\ker J$.
Because $\R^{S\times\Z/p\Z}=V\oplus T$, Lemma~\ref{lem:surjectivity}
implies that $J|_V$ is surjective. The real-analytic implicit-function
theorem therefore makes $(\Phi|_V)^{-1}(0)$ near zero a manifold of dimension
\[
\dim V-m(m-1)=(m-1)(p-m-1),
\]
proving \eqref{eq:dimension}. All hypotheses hold: $\Phi$ is a finite sum
of exponentials of linear functions, $\Phi(0)=0$, and its slice derivative
has just been proved surjective.

If this dimension is positive, take any nonzero coordinate line in a
local analytic chart. It gives an analytic curve $\theta(t)\in V$ of
zeros with nonzero derivative. The resulting matrices $H(\theta(t))$
are not obtained by row and column multiplication of $F_{S,p}$:
the fixed base row and column would force every multiplier to cancel,
making the entire matrix equal to $F_{S,p}$. This argument allows arbitrary
parameter-dependent multipliers; it does not differentiate them.
The finite permutation group cannot turn a positive-dimensional
slice into an isolated point.
If the dimension is zero, the implicit-function chart contains just one
point, so the truncation is locally isolated. For $m=1$ the slice is
already the zero space. This completes Theorem~\ref{thm:dimension}.

\begin{corollary}[Sharp isolation and minimal defect]\label{cor:classification}
A nonempty $m$-row truncation of $F_p$ is isolated exactly for
$m\in\{1,p-1,p\}$, and exactly when its undephased defect is $m+p-1$.
For $m<p$, its undephased defect is $m(p-m+1)$; for $m=p$, it is $2p-1$.
\end{corollary}
\begin{proof}
For proper sets the defect is $mp-m(m-1)$ by Lemma~\ref{lem:surjectivity}.
Theorem~\ref{thm:dimension} gives isolation precisely when the quotient
dimension vanishes. The full prime Fourier matrix is isolated as well
\cite[Section~5.1(i)]{TZ}, \cite[Proposition~1.12]{BOP}.
For a direct proof in the same coordinates,
when $S=\Z/p\Z$, the kernel equations at each nonzero frequency say that
$x_s(d)$ is constant along the whole step-$d$ cycle, hence independent
of $s$. Fourier inversion then gives $A_{sj}=a_s+b_j$.
For $p=2$, the same statement uses one real coefficient at $d=1$.
Thus $\ker J=T$ and $J|_V$ is injective. Select a linear map from its
target to $V$ whose composition with $J|_V$ is the identity. Composing
$\Phi|_V$ with that map gives a square map with invertible derivative;
the inverse-function theorem makes its zero locally unique. In particular
the zero of $\Phi|_V$ is unique. Also $\ker J=T$ has dimension $2p-1$.
\end{proof}

\begin{corollary}[Every tangent is integrable]\label{cor:tangents}
For a proper prime Fourier row set, every $A\in\ker J$ is the phase
velocity at zero of a real-analytic curve of partial Hadamard matrices.
\end{corollary}
\begin{proof}
The surjective derivative makes $\Phi^{-1}(0)$ an analytic manifold
with tangent space $\ker J$. Choose an analytic chart $\psi$ at zero.
Its derivative maps the coordinate space isomorphically onto $\ker J$;
choose $u$ with $D\psi(0)u=A$ and set $\theta(t)=\psi(tu)$.
Then $\theta'(0)=A$ and $H(\theta(t))$ is the required curve.
\end{proof}

\section{An explicit defect formula beyond prime order}
Let $G$ be a finite abelian group, $N=|G|\geq2$, and
$\varnothing\ne S\subseteq G$, $m=|S|$. Index the columns by the
character group $\widehat G$ and put $F_{S,G}(s,\chi)=\chi(s)$.
For $d\ne0$, let $c_d(S)$ count the cosets of $\langle d\rangle$
contained entirely in $S$, and put $C(S)=\sum_{d\ne0}c_d(S)$.
The undephased defect $d(F_{S,G})$ means the real dimension of the
kernel of the orthogonality differential, as in
\cite[Definition~1.4 and Proposition~1.5]{BOP}.

\begin{corollary}[Finite abelian defect and regularity]\label{cor:abelian}
Every such truncation satisfies
\begin{equation}\label{eq:generaldefect}
d(F_{S,G})=m(N-m+1)+C(S).
\end{equation}
Its orthogonality differential is surjective if and only if $C(S)=0$.
In that case its local dephased space is an analytic manifold of
dimension $(m-1)(N-m-1)$.
\end{corollary}
\begin{proof}
Fourier coefficients $x_s(d)=\sum_{\chi\in\widehat G}\chi(d)A_{s\chi}$
give a real-linear bijection between real phase arrays and arrays with
$x_s(-d)=\overline{x_s(d)}$, by character orthogonality and inversion.
The kernel equations are $x_s(d)=x_{s-d}(d)$ whenever $s,s-d\in S$.
Let $k_d$ be the number of connected components of the graph on $S$
joining $s$ to $s+d$ when both are retained. At $d=0$ there are $m$
free real coefficients. For $d\ne-d$, the pair $d,-d$ contributes
$2k_d$ real dimensions; for $d=-d\ne0$, it contributes $k_d$ real
dimensions. Thus $d(F_{S,G})=m+\sum_{d\ne0}k_d$.

Put $n_d=|\{s\in S:s+d\in S\}|$. If the order of $d$ is at least
three, the full step-$d$ graph is a union of cycles and its induced graph
is a union of paths and $c_d(S)$ complete cycles. Hence
$k_d=m-n_d+c_d(S)$. If $d$ has order two, each retained edge is a
whole coset, counted twice in $n_d$; then $n_d=2c_d(S)$ and
$k_d=m-c_d(S)=m-n_d+c_d(S)$ again.
Finally $\sum_{d\ne0}n_d=m(m-1)$ counts the ordered distinct row pairs,
which proves \eqref{eq:generaldefect}. Rank--nullity gives
$\operatorname{rank}_{\R}J=m(m-1)-C(S)$, so regularity is equivalent
to $C(S)=0$. The dephasing and analytic level-set argument of Section~3
then applies unchanged with $p$ replaced by $N$.
\end{proof}

For a proper subset of a prime cyclic group, $c_d(S)=0$ for every
$d\ne0$, recovering Theorem~\ref{thm:dimension}.
For $S=G$, $c_d(G)=[G:\langle d\rangle]$, recovering the known
square formula \cite[Theorem~4.2]{Defect}.
At composite cyclic order the correction is concrete: for
$G=\Z/4\Z$ and $S=\{0,2\}$, only $c_2(S)=1$, so the defect is $7$,
not the prime-case value $6$. For $G=\Z/8\Z$ and $S=\{0,1,2\}$,
no complete nontrivial cyclic coset is retained, so the local dephased
dimension is $(3-1)(8-3-1)=8$.
When $C(S)>0$, formula~\eqref{eq:generaldefect} computes a tangent
dimension; it makes no claim that all tangents integrate.

\section{Coverage and verification}
The later exposition \cite[Conjecture 8.36, pp.~199--200]{Banica}
repeats the density assertion and the two boundary suggestions.
Its Exercise 8.39 asks for defects of Fourier truncations.
Corollary~\ref{cor:abelian} answers that exercise for arbitrary finite
abelian groups and all nonempty row sets. The defect counts first-order
solutions; Theorem~\ref{thm:dimension} and the regular case of
Corollary~\ref{cor:abelian} also integrate them after dephasing.

\begin{center}\small
\begin{tabular}{@{}p{.43\linewidth}p{.50\linewidth}@{}}\toprule
Source item & Mathematical coverage\\\midrule
\cite{BOP}, Conj.~1.14; \cite{Banica}, Conj.~8.36 &
Refuted-and-corrected: isolation iff $m\in\{1,p-1,p\}$, by
Corollaries~\ref{cor:counterexample} and \ref{cor:classification}.\\
Paragraph after Conj.~1.14: $p-2$ rows &
Refuted-and-corrected: non-isolated for every prime $p\geq5$,
dimension $p-3$.\\
Same paragraph: $p-1$ rows & Proved: isolated by the zero-dimensional slice.\\
\cite{BOP}, Prop.~1.12 & Prior proof cited: full-prime rigidity;
also proved directly in Corollary~\ref{cor:classification}.\\
\cite{BOP}, Conj.~1.8; \cite{Banica}, Conj.~8.32 &
Proved on the prime Fourier family, Corollary~\ref{cor:classification};
the general conjecture is not covered.\\
Comment after \cite{BOP}, Conj.~1.8 (p.~7) &
Not posed for this family; proved: every proper-prime tangent
integrates by Corollary~\ref{cor:tangents}. The suggested broader
phenomenon is not covered.\\
\cite{BOP}, Conj.~4.5 & Not covered for its output matrices.
Its isolated Fourier inputs have exactly $m=1,p-1,p$;
the stated input hypothesis is retained.\\
\cite{BOP}, Conjs.~2.4--2.6, 2.10, 3.7, 3.9, 6.3--6.4 &
Not covered in their general settings: regularity, master matrices,
and almost Hadamard matrices require additional arguments.\\
\cite{Banica}, Ex.~8.39 & Answered: yes. Formula~\eqref{eq:generaldefect}
gives the defect for every finite abelian group and row set.\\
General regularity criterion & Not posed; proved:
$C(S)=0$ exactly characterizes surjectivity, with analytic dimension
$(m-1)(N-m-1)$. Singular composite isolation is not covered.\\\bottomrule
\end{tabular}
\end{center}

The semigroup classification question after \cite[Theorem~5.9, p.~22]{BOP}
and the quantum-deformation question after \cite[Definition~6.2, p.~22]{BOP}
are not covered. The square master classification of
\cite[Conjecture~3.9]{BOP} remains a square question: the isolated
rectangular master matrices $F_{p-1,p}$ do not refute it.

The written proofs above establish the complete dimension and isolation
classification. The following declarations, each linked to its exact
Lean source, give the indicated formal statements.
\begin{itemize}
\item \href{https://github.com/the-omega-institute/trureturing/blob/43ced8a8460ce89a7c5cae3deeac10e5fc08d565/D5/S3/Quantum/Tomography/Hadamard/BanicaCycleForestIncidence.lean}{\texttt{BanicaCycleForestIncidence.lean}}:
\par\texttt{cycle\_incidence\_surjective} integrates edge data on a proper
subset with unit step, for any positive modulus and abelian coefficient group.
\item \href{https://github.com/the-omega-institute/trureturing/blob/43ced8a8460ce89a7c5cae3deeac10e5fc08d565/D5/S3/Quantum/Tomography/Hadamard/BanicaTruncationJacobian.lean}{\texttt{BanicaTruncationJacobian.lean}}:
\texttt{fourier\_jacobian\_surjective} proves the prime proper-subset
surjectivity, and \texttt{phase\_jacobian\_rank\_nullity} gives the kernel count.
\item \href{https://github.com/the-omega-institute/trureturing/blob/43ced8a8460ce89a7c5cae3deeac10e5fc08d565/D5/S3/Quantum/Tomography/Hadamard/BanicaLevelSetCurve.lean}{\texttt{BanicaLevelSetCurve.lean}}:
\texttt{analytic\_level\_set\_curve} realizes a regular analytic
level-set tangent by a global real-analytic curve, up to a positive scalar.
Its regularity parameter \texttt{ContDiff} $\R$ $\top$ denotes
analytic order $\omega$ in the source's Mathlib version.
\item \href{https://github.com/the-omega-institute/trureturing/blob/43ced8a8460ce89a7c5cae3deeac10e5fc08d565/D5/S3/Quantum/Tomography/Hadamard/BanicaTruncationRefutation.lean}{\texttt{BanicaTruncationRefutation.lean}}:
\texttt{two\_row\_truncation\_not\_isolated} proves non-isolation for every
prime $p\geq5$ and every $|S|=p-2$; \texttt{result} proves the negation
of the uniform-density assertion. Its trivial class permits arbitrary
row and column phase functions, and so includes the fixed-exponent
trivial curves of \cite[Definition 1.3]{BOP}.
\end{itemize}
The sharp slice dimension, tangent integration with an exact local
velocity, one-row boundary, isolation classification and finite abelian
formula are written proofs, not the statements of the final Lean result.
The accompanying \texttt{verify.py} and \texttt{verify\_cyclic.py} give
exact finite-field Jacobian ranks and finite graph checks for the
explicitly listed cyclic and finite abelian examples. These checks
illustrate the formulas rather than replace their proofs.

\begin{thebibliography}{9}
\bibitem{TZ} W.~Tadej and K.~\.{Z}yczkowski, \emph{Defect of a unitary matrix},
Linear Algebra Appl. \textbf{429} (2008), 447--481.
\href{https://doi.org/10.1016/j.laa.2008.02.036}{doi:10.1016/j.laa.2008.02.036}.
Text locators refer to \href{https://arxiv.org/abs/math/0702510v2}{arXiv:math/0702510v2}
(20 December 2007).
\bibitem{Defect} T.~Banica, \emph{The defect of generalized Fourier matrices},
Linear Algebra Appl. \textbf{438} (2013), 3667--3688.
\href{https://doi.org/10.1016/j.laa.2013.01.011}{doi:10.1016/j.laa.2013.01.011}.
Text locators refer to \href{https://arxiv.org/abs/1210.2556v2}{arXiv:1210.2556v2}
(8 January 2013).
\bibitem{BOP} T.~Banica, D.~\"Ozteke and L.~Pittau,
\emph{Isolated partial Hadamard matrices, and related topics},
Open Systems \& Information Dynamics \textbf{25} (2018), no.~2, 1850009.
\href{https://doi.org/10.1142/S1230161218500099}{doi:10.1142/S1230161218500099}.
Text locators refer to \href{https://arxiv.org/abs/1706.00986v2}{arXiv:1706.00986v2}
(7 October 2017).
\bibitem{Banica} T.~Banica, \emph{Invitation to Hadamard matrices},
\href{https://arxiv.org/abs/1910.06911v9}{arXiv:1910.06911v9}
(29 July 2024), Section~8d, pp.~195--200.
\end{thebibliography}
\noindent\small\textcopyright\ 2026 Haobo Ma and Wenlin Zhang.
Licensed under Creative Commons Attribution 4.0 International
(\href{https://creativecommons.org/licenses/by/4.0/}{CC BY 4.0}).
\end{document}
